\documentclass{beamer}
\usetheme{Madrid}
\usecolortheme{default}
\usepackage{fontspec}
\usepackage{polyglossia}
\setdefaultlanguage{english}
\usepackage{tikz}
\usepackage{hyperref}
\usepackage{graphicx}
\usepackage{booktabs}

% Custom colors
\definecolor{RurDynamixGreen}{RGB}{76, 175, 80}
\definecolor{RurDynamixBlue}{RGB}{33, 150, 243}
\definecolor{DarkGray}{RGB}{64, 64, 64}

\setbeamercolor{primary}{fg=white,bg=RurDynamixGreen}
\setbeamercolor{secondary}{fg=white,bg=RurDynamixBlue}
\setbeamertemplate{frametitle}{
  \begin{beamercolorbox}[rounded=true,shadow=true,leftskip=1cm,rightskip=1cm,wd=\paperwidth,ht=0.8cm,dp=0.4cm]{frametitle}
    \usebeamerfont{frametitle}\insertframetitle
  \end{beamercolorbox}
}

% Title Page Setup
\title{\tiny{Rural Territories in Transition: Morocco and the Global South}}
\author{Said Moukrim}
\institute{UM5-FSR Morocco}
\date{November 28, 2026}

\begin{document}

\begin{frame}

\begin{columns}[T]
  \column{0.5\textwidth}
  \includegraphics[width=\textwidth]{images/rur_dyn.png}

  \column{0.25\textwidth}
  \includegraphics[width=\textwidth]{images/FSR logo.png}
\end{columns}

\vspace{1cm}

\begin{center}
\Large {Hydraulic Innovation and Rural Transitions:}\\
\vspace{1cm}
\large{Spatio-Temporal Analysis of Artificial Basins Dynamics in Bab Berred}\\
Central Rif (Morocco)\\
\vspace{1.5cm}
\Large{Said Moukrim}
\end{center}

\end{frame}

% Slide 2: Problem Statement
\begin{frame}
\frametitle{The Challenge: Water Scarcity in Mountain Agriculture}
\begin{itemize}
  \item \textbf{Mediterranean mountain climate}: Irregular precipitation, seasonal drought stress
  \item \textbf{Regional trend}: Declining precipitation (Mann-Kendall test, tau = -0.1376, p < 0.05)
  \item \textbf{2023 record}: Lowest precipitation on record (\textasciitilde 300 mm)
  \item \textbf{Agricultural pressure}: Cannabis varieties require 4-5× more water than traditional kif
  \item \textbf{Farmer response}: Massive artificial basin construction (2010-2025)
\end{itemize}
\end{frame}

% Slide 3: Why Cannabis Transformation?
\begin{frame}
\frametitle{The Driver: Economic Opportunity, Not Legalization}
\textbf{Key distinction:}
\begin{itemize}
  \item \textbf{Legalization (2021)}: Law 13-21 enabled official channels
  \item \textbf{BUT} actual driver = international market demand for high-quality cannabis
  \item \textbf{Productivity gap}: Local kif $\approx$ 15 kg/ha $\to$ Modern hybrids $\approx$ 70 kg/ha
  \item \textbf{This 4-5× increase} explains rapid farmer adoption of hybrid varieties
  \item \textbf{Actual cultivated area} $>>$ officially authorized area (gap persists post-2021)
\end{itemize}

\textbf{Conclusion}: International demand, not legalization, is the primary driver
\end{frame}

% Slide 4: Study Area - Bab Berred
\begin{frame}
\frametitle{Study Area: Bab Berred Commune}
\begin{columns}[T]
  \column{0.5\textwidth}
  \textbf{Geographic profile:}
  \begin{itemize}
    \item Location: Chefchaouen Province, Central Rif
    \item Area: 17,224 hectares (172.24 km²)
    \item Altitude: 650-2,100 m
    \item Population: 24,742 residents
    \item Slope: 80% average gradient (mountainous)
  \end{itemize}

  \column{0.5\textwidth}
  \textbf{Climatic context:}
  \begin{itemize}
    \item Mean annual precipitation: $>$ 700 mm
    \item Concentration: Winter (Oct-Apr)
    \item Summer: Hot, dry (May-Sept)
    \item Mean annual temperature: 18.6°C
    \item June mean: 24.6°C
  \end{itemize}
\end{columns}

\begin{center}
  \small{\textit{Figure 1: Location and topography of Bab Berred}}
\end{center}
\end{frame}

% Slide 5: Basin Proliferation - The Numbers
\begin{frame}
\frametitle{Results: Artificial Basin Explosion in Bab Berred}

\textbf{Basin inventory (2010-2025):}
\begin{itemize}
  \item 2010: $\approx$ 28 basins
  \item 2025: 2,044 basins
  \item Overall density: 11.87 basins/km² (mean spacing: $\approx$ 290 m)
  \item Growth pattern: Gradual (2010-2014), then accelerating (2015+)
  \item 2024-2025 alone: 372 new basins
\end{itemize}

\textbf{Storage capacity transformation:}
\begin{itemize}
  \item 2010: $\approx$ 50,000 m³ total
  \item 2025: $\approx$ 4.3 million m³ total
  \item Increase: 85-fold expansion in 15 years
\end{itemize}

\begin{center}
  \small{\textit{Figure 2: Basin evolution curve (2010-2025)}}
\end{center}
\end{frame}

% Slide 6: Basin Characteristics - Miniaturization Pattern
\begin{frame}
\frametitle{Results: Basin Morphometrics and Miniaturization}

\textbf{Size trend:}
\begin{itemize}
  \item 2010 mean basin: 1,260 m²
  \item 2025 mean basin: 812 m²
  \item 2025: 90\% of basins $<$ 1,500 m²
\end{itemize}

\textbf{Depth and volume:}
\begin{itemize}
  \item Mean depth: 5.5 m
  \item Mean storage volume per basin: $\approx$ 2,100 m³
  \item Pattern: Multiple small reservoirs (economically constrained) rather than fewer large structures
\end{itemize}

\textbf{Interpretation:}
\begin{itemize}
  \item Reflects resource constraints of rural communities
  \item Incremental infrastructure investment
  \item Pragmatic adaptation to limited capital
\end{itemize}

\begin{center}
  \small{\textit{Figure 3: Basin size distribution over time}}
\end{center}
\end{frame}

% Slide 7: Spatial Distribution Factors
\begin{frame}
\frametitle{Results: Where Are Basins Located?}

\textbf{Altitude control:}
\begin{itemize}
  \item Concentration: 800-1,500 m elevation
  \item Reflects balance between water demand and source proximity
\end{itemize}

\textbf{Slope preference:}
\begin{itemize}
  \item 71\% of basins on slopes 5-15\%
  \item Virtually absent on slopes $>$ 25\%
  \item Reflects engineering constraints
\end{itemize}

\textbf{Land-use distribution:}
\begin{itemize}
  \item 78\% within cultivated lands
  \item 18\% in grasslands
  \item Only 3\% in natural vegetation
  \item Shows deliberate agricultural intensification + land conversion
\end{itemize}

\begin{center}
  \small{\textit{Figures 4-5: Maps of altitude, slope, and land-use distribution}}
\end{center}
\end{frame}

% Slide 8: Precipitation Decline
\begin{frame}
\frametitle{The Paradox Part 1: Declining Precipitation (1985-2025)}

\textbf{Long-term trend:}
\begin{itemize}
  \item Progressive decline in annual precipitation over 40 years
  \item Mann-Kendall test: tau = -0.1376 (significant negative trend, p $<$ 0.05)
  \item High interannual variability masks overall decline
\end{itemize}

\textbf{Extreme variability:}
\begin{itemize}
  \item Wet years: 1996, 2010, 2018 ($>$ 850 mm)
  \item Dry years: 1995, 1998, 2005, 2015, 2020 ($<$ 410 mm)
  \item Extreme: 2023 record low at $\approx$ 300 mm
\end{itemize}

\textbf{Agricultural consequence:}
\begin{itemize}
  \item Increasing hydrological pressure on rainfed agriculture
  \item Motivation for artificial water storage as drought mitigation
\end{itemize}

\begin{center}
  \small{\textit{Figure 6: Precipitation time series and trend (1985-2025)}}
\end{center}
\end{frame}

% Slide 9: NDVI Paradox
\begin{frame}
\frametitle{The Paradox Part 2: Vegetation Improvement Despite Drought}

\textbf{NDVI analysis (March-July, 1998-2025):}
\begin{itemize}
  \item 1998-2012: Slight increase (Sen's slope 0.0025 per year)
  \item 2012 onward: Pronounced acceleration
  \item Homogeneity break detected 2012 (p = 0.00000434, tau = 0.501)
  \item Peak amplitudes: Highest during cannabis growing season (May-Aug)
\end{itemize}

\textbf{The paradox:}
\begin{itemize}
  \item Precipitation is declining overall
  \item NDVI is improving strongly
  \item Timing: Both shifts occur around 2012-2013
  \item Conclusion: Irrigation infrastructure compensates for natural water deficit
\end{itemize}

\begin{center}
  \small{\textit{Figure 7a: NDVI trends (1998-2025)}}
\end{center}
\end{frame}

% Slide 10: What NDVI Really Shows
\begin{frame}
\frametitle{The Deception: What NDVI Improvement Really Represents}

\textbf{What NDVI indicates:}
\begin{itemize}
  \item Enhanced biomass productivity
  \item NOT ecosystem health or biodiversity improvement
\end{itemize}

\textbf{The truth behind the signal:}
\begin{itemize}
  \item NDVI improvement derives from irrigated cannabis monoculture
  \item Cannabis NDVI values $>$ 0.75 (among highest globally)
  \item Peak growth March-July aligns exactly with measurement period
  \item High biomass signal masks ecological degradation
\end{itemize}

\textbf{What actually happened:}
\begin{itemize}
  \item Natural areas (grassland, rangeland) $\to$ declined
  \item Cultivated agricultural area $\to$ expanded
  \item Land-use transformation: Natural ecosystems $\to$ intensive cannabis
  \item NOT ecosystem-wide improvement; ecosystem conversion!
\end{itemize}

\begin{center}
  \small{\textit{Figure 8: Land-use change map comparison}}
\end{center}
\end{frame}

% Slide 11: Hydrological Sustainability Crisis
\begin{frame}
\frametitle{The Critical Concern: Aquifer Overexploitation}

\textbf{The hydrological crisis:}
\begin{itemize}
  \item Water extraction $>$ annual renewable recharge capacity
  \item Progressive piezometric (water table) decline occurring
  \item Threatens long-term water security
\end{itemize}

\textbf{Why this matters:}
\begin{itemize}
  \item Incompatible with hydrological system recovery
  \item Precipitation deficits intensifying with climate change
  \item Cumulative impacts of 2,000+ small reservoirs create complex pressure
  \item May trigger unintended ecosystem consequences
\end{itemize}

\textbf{The paradox of adaptation:}
\begin{itemize}
  \item Infrastructure enables short-term productivity
  \item But may cause long-term water insecurity
  \item Rational individual decisions (basin construction) lead to collectively unsustainable outcome (aquifer depletion)
\end{itemize}

\begin{center}
  \small{\textit{Figure 9: Aquifer piezometry decline (conceptual)}}
\end{center}
\end{frame}

% Slide 12: Ecological Consequences
\begin{frame}
\frametitle{Beyond Hydrology: Landscape Fragmentation and Biodiversity Loss}

\textbf{Landscape fragmentation:}
\begin{itemize}
  \item Natural spaces increasingly fragmented by basin infrastructure
  \item Cannabis cultivation expanding through grasslands
  \item Disrupts wildlife habitat connectivity
  \item Threatens endemic Mediterranean species in restricted refugia
\end{itemize}

\textbf{Biodiversity threat in Rif Mountains:}
\begin{itemize}
  \item Rif harbors important endemic plant species
  \item Several species in restricted refugia within/adjacent to Bab Berred
  \item Habitat fragmentation directly threatens species persistence
  \item Ecosystem functionality compromised
\end{itemize}

\textbf{Ecological summary:}
\begin{itemize}
  \item Apparent NDVI improvement masks simultaneous ecological deterioration
  \item Natural area regression + cultivated area expansion
  \item Remote sensing alone insufficient for sustainability assessment
\end{itemize}

\begin{center}
  \small{\textit{Figure 10: Landscape transformation 2006 vs 2023}}
\end{center}
\end{frame}

% Slide 13: Governance Implications
\begin{frame}
\frametitle{Governance and Sustainable Futures}

\textbf{The current problem:}
\begin{itemize}
  \item Absence of coordinated management
  \item Individual cultivators optimize household advantage
  \item Individually rational $\neq$ collectively sustainable
  \item Results: Aquifer depletion, landscape fragmentation, biodiversity loss
\end{itemize}

\textbf{Necessary reforms:}
\begin{itemize}
  \item Hydrological assessments imposing sustainable extraction ceilings
  \item Protection of natural habitats
  \item Agricultural diversification
  \item Improved water efficiency technologies
  \item Land-use planning integrating conservation with agriculture
\end{itemize}

\textbf{Role of remote sensing:}
\begin{itemize}
  \item NDVI alone suggests environmental improvement
  \item Integrated analysis (NDVI + precipitation + land-use) reveals true challenges
  \item Essential for evidence-based governance
\end{itemize}

\begin{center}
  \small{\textit{Applies globally to similar Mediterranean regions}}
\end{center}
\end{frame}

% Slide 14: Conclusion
\begin{frame}
\frametitle{Conclusion: Rural Adaptation and Global Relevance}

\textbf{What we found:}
\begin{itemize}
  \item Rapid artificial basin proliferation (28 $\to$ 2,044 in 15 years)
  \item Infrastructure compensates for declining precipitation
  \item Apparent vegetation improvement masks ecological degradation
  \item Aquifer overexploitation threatens long-term sustainability
\end{itemize}

\textbf{Why it matters beyond Bab Berred:}
\begin{itemize}
  \item Pattern replicated across Mediterranean and semi-arid regions globally
  \item Infrastructure-intensive agricultural transformation widespread
  \item Remote sensing-based monitoring often misses underlying ecological decline
  \item Governance frameworks must transcend infrastructure investment alone
\end{itemize}

\textbf{Path forward:}
\begin{itemize}
  \item Sustainable rural futures require integrated governance
  \item Simultaneous attention to: hydrological limits, ecosystem integrity, livelihood resilience
  \item Particularly urgent as climate change intensifies water scarcity
\end{itemize}

\end{frame}

% Final Slide
\begin{frame}[plain]
\begin{center}
  \vfill
  {\Huge \textbf{Thank you for your attention}}
  \\[1cm]
  {\Large Questions?}
  \\[2cm]
  \small{
    Ousmane Samba MAÏGA, Imane Aouassar, et al. \\
    Mohammed-V University in Rabat \\
    Faculty of Sciences, Laboratory of Botany, Mycology and Environment
  }
  \vfill
\end{center}
\end{frame}

\end{document}
